\documentclass[../../main]{subfiles}

\begin{document}

\section{Lógica de Predicados}
\label{sec:matematica:logica-matematica:logica-de-predicados}

\begin{defn}[Predicado]
	\label{def:matematica:logica-matematica:logica-de-predicados:predicado}

	Um predicado é uma expressão que representa uma propriedade ou uma
	relação envolvendo uma ou mais variáveis.

	Quando valores são atribuídos às variáveis, o predicado transforma-se
	em uma proposição, podendo então ser classificado como verdadeiro ou
	falso.
\end{defn}

\begin{ex}

	Considere o predicado

	\[
		P(x): x > 0.
	\]

	Se atribuirmos \(x=3\), obtemos a proposição

	\[
		P(3): 3 > 0,
	\]

	que é verdadeira.

\end{ex}

\begin{defn}[Domínio de Discurso]
	\label{def:matematica:logica-matematica:logica-de-predicados:dominio-de-discurso}

	O domínio de discurso é o conjunto dos objetos sobre os quais as
	variáveis de um predicado podem assumir valores.
\end{defn}

\begin{defn}[Variável Livre]
	\label{def:matematica:logica-matematica:logica-de-predicados:variavel-livre}

	Uma variável é denominada livre quando não está associada ao escopo de
	nenhum quantificador.
\end{defn}

\begin{defn}[Variável Ligada]
	\label{def:matematica:logica-matematica:logica-de-predicados:variavel-ligada}

	Uma variável é denominada ligada quando aparece no escopo de um
	quantificador.
\end{defn}

\begin{defn}[Quantificador]
	\label{def:matematica:logica-matematica:logica-de-predicados:quantificador}

	Um quantificador é um operador lógico utilizado para especificar sobre
	quantos elementos do domínio um predicado deve ser considerado.
\end{defn}

\begin{defn}[Quantificador Universal]
	\label{def:matematica:logica-matematica:logica-de-predicados:quantificador-universal}

	Seja \(P(x)\) um predicado definido sobre um domínio \(D\).

	A expressão

	\[
		\forall x \, P(x)
	\]

	afirma que \(P(x)\) é verdadeira para todo elemento
	\(x \in D\).
\end{defn}

\begin{defn}[Quantificador Existencial]
	\label{def:matematica:logica-matematica:logica-de-predicados:quantificador-existencial}

	Seja \(P(x)\) um predicado definido sobre um domínio \(D\).

	A expressão

	\[
		\exists x \, P(x)
	\]

	afirma que existe pelo menos um elemento
	\(x \in D\) para o qual \(P(x)\) é verdadeira.
\end{defn}

\begin{prop}[Negação do Quantificador Universal]
	\label{prop:matematica:logica-matematica:logica-de-predicados:negacao-quantificador-universal}

	Para todo predicado \(P(x)\),

	\[
		\neg(\forall x\,P(x))
		\Leftrightarrow
		\exists x\,\neg P(x).
	\]

\end{prop}

\begin{proof}

	A afirmação

	\[
		\forall x\,P(x)
	\]

	significa que todos os elementos do domínio satisfazem \(P\).

	Sua negação afirma que existe pelo menos um elemento que não satisfaz
	\(P\).

	Portanto,

	\[
		\neg(\forall x\,P(x))
		\Leftrightarrow
		\exists x\,\neg P(x).
	\]

\end{proof}

\begin{prop}[Negação do Quantificador Existencial]
	\label{prop:matematica:logica-matematica:logica-de-predicados:negacao-quantificador-existencial}

	Para todo predicado \(P(x)\),

	\[
		\neg(\exists x\,P(x))
		\Leftrightarrow
		\forall x\,\neg P(x).
	\]

\end{prop}

\begin{proof}

	A afirmação

	\[
		\exists x\,P(x)
	\]

	significa que existe ao menos um elemento do domínio para o qual
	\(P\) é verdadeira.

	Sua negação afirma que nenhum elemento satisfaz \(P\).

	Portanto,

	\[
		\neg(\exists x\,P(x))
		\Leftrightarrow
		\forall x\,\neg P(x).
	\]

\end{proof}

\begin{prop}[Distribuição do Quantificador Universal sobre a Conjunção]
	\label{prop:matematica:logica-matematica:logica-de-predicados:universal-conjuncao}

	Para quaisquer predicados \(P(x)\) e \(Q(x)\),

	\[
		\forall x \bigl(P(x)\land Q(x)\bigr)
		\Leftrightarrow
		(\forall x\,P(x))
		\land
		(\forall x\,Q(x)).
	\]

\end{prop}

\begin{proof}

	A proposição da esquerda afirma que todo elemento do domínio satisfaz
	simultaneamente \(P\) e \(Q\).

	Isso ocorre exatamente quando todo elemento satisfaz \(P\) e todo
	elemento satisfaz \(Q\).

	Portanto as duas expressões são logicamente equivalentes.

\end{proof}

\begin{prop}[Distribuição do Quantificador Existencial sobre a Disjunção]
	\label{prop:matematica:logica-matematica:logica-de-predicados:existencial-disjuncao}

	Para quaisquer predicados \(P(x)\) e \(Q(x)\),

	\[
		\exists x \bigl(P(x)\lor Q(x)\bigr)
		\Leftrightarrow
		(\exists x\,P(x))
		\lor
		(\exists x\,Q(x)).
	\]

\end{prop}

\begin{proof}

	A expressão da esquerda afirma que existe um elemento que satisfaz
	\(P\) ou \(Q\).

	Isso ocorre exatamente quando existe um elemento que satisfaz \(P\) ou
	existe um elemento que satisfaz \(Q\).

	Portanto as duas expressões são equivalentes.

\end{proof}

\end{document}
